\section{Zadanie 7}
Niech \(T_n\) będzie pełnym drzewem binarnym na \(n\) wierzchołkach. Rozważmy spacer losowy po \(T_n\), gdzie z prawdopodobieństwem \(\frac{1}{2}\) stoimy w miejscu, a jeśli mamy wykonać ruch to wykonujemy go jednostajnie na jednego z sąsiadów. Zauważ, że rozkład stacjonarny odpowiednio zdefiniowanego łańcucha Markova istnieje. Pokaż, że \(\tau_{mix}\) jest liniowy w terminach \(n\).\\
\\
\textit{Wskazówka:} Jeśli \(p,q,r \in [0,1]\) takie, że \(p+q+r=1\), \(p \neq q\) i rozważamy spacer na liczbach \([0,n]\) taki, że prawdopodobieństwo pozostania w miejscu to \(r\), prawdopodobieństwo przejścia z \(i\) do \(i-1\) wynosi \(p\), a prawdopodobieństwo przejścia z \(i\) do \(i+1\) wynosi q. Jeśli jesteśmy w punkcie 0 i mamy iść w lewo to zostajemy w miejscu, a jeśli jesteśmy w punkcie \(n\) i mamy iść w prawo to również zostajemy w miejscu. Niech \(X\) będzie zmienną losową będącą liczbą kroków potrzebnych do przejścia z \(n\) do \(0\). Wtedy
\[
	E(X) \leq \frac{1}{p-q}\left[n-q\left(\frac{1-(q/p)^n}{p-q}\right)\right].
\]

\subsection*{Rozwiązanie}

Zacznijmy od rozkładu stacjonarnego. Łańcuch jest skończony i nieprzywiedlny, więc taki rozkład istnieje. \\

Rozważmy teraz sprzęganie łańcuchów Markova \(((X_t, Y_t))_{t\in \mathbb{N}}\), których to macierz przejścia jest określona naszym grafem(możemy ponumerować wierzchołki od \(0\) do \(n-1\)). Określamy \(X_{t+1}, Y_{t+1}\) na podstawie \(X_t, Y_t\) w następujący sposób:
\begin{itemize}
	\item Jeśli \(X_t, Y_t\) były wierzchołkami z różną odległością od korzenia drzewa, to z prawdopodobieństwem \(\frac{1}{2} : X_{t+1} = X_t, \) nastomiast \(Y_{t+1}\) wybieramy jednostajnie spośród sąsiadów \(Y_t\), i z prawdopodobieństwem \(\frac{1}{2}:Y_{t+1} = Y_t, \) nastomiast \(X_{t+1}\) wybieramy jednostajnie spośród sąsiadów \(X_t\).
	\item Jeśli \(X_t, Y_t\) na tym samym poziomie drzewa, to wybieramy ten ruch który wykonujemy na obu tych zmiennych - czyli z prawdopodobieństwem \(\frac{1}{2}\) będzie \(X_{t+1} = X_t, Y_{t+1}=Y_t\), albo wybieramy jednostajnie z sąsiadów tych wierzchołków do którego idziemy(w obu idziemy do lewego dziecka/w obu do rodzica/w obu do prawego dziecka). To zapewnia, że gdy raz \(X_t, Y_t\) będą na jednym poziomie w drzewie, to już zawsze będą na tym samym poziomie.
\end{itemize}
Niech \(Z = Z_1+Z_2\) - oczekiwana liczba kroków po których \(X_t = Y_t\)(tu już chodzi nam o konkretnie ten sam wierzchołek, nie tylko poziom). \\
\\
\(Z_1\) to oczekiwana liczba kroków po których \(X_t, Y_t\) znajdą się na tym samym poziomie. Zauważmy, że \(D_t = |\text{poziom}(X_t)-\text{poziom}(Y_t)|\) możemy sprowadzić do spaceru symetrycznego z dwoma barierami w \(0\) i \(\log_2n\), bo w każdym kroku różnica między poziomami zmienia się o co najwyżej 1. Z przypomnienia o spacerach losowych oczekiwana liczba kroków do uderzenia w 0 niezależnie od miejsca startu mieści się w \(O((\log_2n)^2)\).\\
\\
\(Z_2\) to oczekiwana liczba kroków po których \(X_t = Y_t\) od momentu gdy znajdą się na tym samym poziomie w drzewie. Zauważmy, że jest ona nie większa niż oczekiwana liczba kroków przejścia \(X_t, Y_t\) z liści do korzenia, co znowu możemy zinterpretować przez spacer niesymetryczny z dwoma barierami. Czyli ze wskazówki mamy(u nas \(p = \frac{1}{6}, q=\frac{1}{3})\):
\[
	E[Z_2] \leq \frac{1}{p-q}\left[\log_2n-q\left(\frac{1-(q/p)^{\log_2n}}{p-q}\right)\right] = -6\log_2n - 12 + 12 \cdot 2^{\log_2n} = \theta(n)
\]

Wyżej zostało zamienione \(n\) ze wskazówki na \(\log_2n\), bo tutaj jest spacer ograniczony przez 0 i wysokość drzewa, czyli \(\log_2n\).

Mamy więc, że \(E[Z] \leq cn\) dla pewnego \(c \in \mathbb{R}\).

Korzystamy teraz z lematu \(10.3.2\), przytoczonego m.in. w zadaniu 2. Tutaj mamy(po drodze korzystając z nierówności Markova):
\[
	P(X_t \neq Y_t | X_0 = x, Y_0 = y) = P(Z > t) \leq \frac{E[Z]}{t} \leq \frac{cn}{t}
\]
Dla \(t = 4cn\) mamy:
\[
	P(X_t \neq Y_t | X_0 = x, Y_0 = y) \leq \frac{1}{4}
\]
Więc z lematu dostajemy
\[
	\tau_{mix} = \tau_{mix}(\epsilon) = \tau_{mix}(\frac{1}{4}) \leq 4cn = O(n)
\]

